%% CV - Giovanny Garcia
%% Tailored for Epic Games, Tools Programmer Intern (R27410)
%% Compile with: lualatex main_epic_games.tex

\documentclass[11pt,a4paper,sans]{moderncv}
\moderncvstyle{banking}
\moderncvcolor{blue}

\renewcommand*{\sectionstyle}[1]{{\sectionfont\color{color1}#1}}
\makeatletter
\AtBeginDocument{%
  \@ifundefined{firstnamestyle}{}{\renewcommand*{\firstnamestyle}[1]{{\fontsize{34}{36}\bfseries\upshape\color{color1}#1}}}%
  \@ifundefined{lastnamestyle}{}{\renewcommand*{\lastnamestyle}[1]{{\fontsize{34}{36}\bfseries\upshape\color{color1}#1}}}%
}
\makeatother

\usepackage[utf8]{inputenc}
\usepackage[scale=0.80]{geometry}
\usepackage{import}
\usepackage{needspace}

% personal data
\name{Giovanny}{Garcia}
\address{Austin, Texas, United States}{}{}
\email{giovanny.garcia2@g.austincc.edu}
\extrainfo{\href{https://linkedin.com/in/giovanny-garcia-07ab24322}{linkedin.com/in/giovanny-garcia-07ab24322}, \href{https://github.com/giovanny-garcia}{github.com/giovanny-garcia}}

\begin{document}
\hypersetup{
    colorlinks=true,
    linkcolor=blue,
    filecolor=magenta,
    urlcolor=blue,
    pdftitle={Giovanny Garcia - CV},
    pdfpagemode=FullScreen,
}

\makecvtitle

% ============================================================
%     PROFILE STATEMENT
% ============================================================

\vspace{6pt}
\small{Computer Science graduate (Associate of Science, Austin Community College) who builds tools other people can run: a TypeScript Discord bot designed around a 3-request/day API quota, and a Claude Code workflow that turns job-search admin into a repeatable process. Seeking Epic's paid Tools Programmer Intern role on the Build Operations Tools team (May 3 to August 27, 2027) to help developers write, build, test, and ship Fortnite and Unreal Engine work with fewer defects and less friction.}

% ============================================================
%     CORE COMPETENCIES
% ============================================================

\section{Core Competencies}
\vspace{1pt}
\begin{itemize}

\item \textbf{Agentic workflows and AI-assisted coding:} Use Claude Code to evaluate messy source material, generate structured documents, and keep every claim tied to a file. Comfortable directing an AI coding agent, rejecting bad output, and iterating until someone else can rerun the workflow.

\item \textbf{Developer tools and command UX:} Ship slash-command tools, embed layouts, and settings flows in discord.js so a small group can browse data, subscribe to alerts, and check quota without opening the vendor site. Not a desktop GUI-framework portfolio. Strong on interaction design for tools people actually click.

\item \textbf{Data processing, caching, and usage reporting:} REST clients, SQLite schemas, upserts, and a cache-first sync so daily use does not burn a tight API budget. A \texttt{/quota} command reports daily usage and cache status, which is the same instinct as instrumenting a tool so the team can see how it is used.

\item \textbf{Automation, packaging, and reliability:} TypeScript compile to \texttt{dist/}, Node 18+ runtime, poll loops that read cache only, and exponential backoff capped at one hour when an announce pass fails. README-level setup so another person can run the system without guessing.

\item \textbf{Languages for this posting:} Production work in TypeScript and JavaScript. Coursework foundation in C++ and C\# (LinkedIn-listed; no public Unreal or C\# game repo yet). Ready to learn C++ on the job, which the posting explicitly invites.

\end{itemize}

% ============================================================
%     INDEPENDENT PROJECTS
% ============================================================

\section{Independent Projects}
\vspace{3pt}
\begin{itemize}

\needspace{5\baselineskip}
\item{\cventry{2026}{CS2 Discord Bot (TypeScript, Node.js, SQLite)}{Personal project}{Austin, TX}{}{\vspace{1pt}
\href{https://github.com/giovanny-garcia/hltv-discord-bot}{github.com/giovanny-garcia/hltv-discord-bot}
\begin{itemize}
    \item Built a Discord bot that turns public Counter-Strike 2 match data (GGScore API) into slash commands, embed announcements, and starting-soon reminders a small group can use without checking the source by hand.
    \item Designed a cache-first architecture: API calls happen only on an explicit \texttt{/sync}. Browse commands, announcements, and the poll loop read SQLite so the 3-request/day free tier is not wasted on routine use.
    \item Split the code into services (API client, cache, announce, poll), command modules, SQLite storage, and embed helpers. Persistence covers guild settings, seen-item deduplication, cached datasets, and daily API usage counters.
    \item Implemented a poll loop that posts new-match and starting-soon reminders from cache, with exponential backoff (capped at one hour) when a pass fails.
    \item Designed the operator UX: \texttt{/subscribe}, \texttt{/settings}, and \texttt{/quota} (ephemeral usage and cache status). Manage Server permission gates sync and settings. Documented the free-tier call budget so another person can run the bot without burning quota on boot.
    \item Packaged the project for someone else to run: Node 18+, \texttt{tsx} for development, \texttt{tsc} to \texttt{dist/}, \texttt{.env.example}, and a README covering setup, environment variables, and a recommended first sync.
\end{itemize}}}

\vspace{3pt}

\needspace{5\baselineskip}
\item{\cventry{2026}{Applied AI job-search workflow (Claude Code)}{Personal project}{Austin, TX}{}{\vspace{1pt}
\href{https://github.com/giovanny-garcia/ai-job-search}{github.com/giovanny-garcia/ai-job-search}
\begin{itemize}
    \item Forked and ran an open-source application framework on Claude Code to evaluate job postings, tailor a CV, and draft cover letters from a structured profile.
    \item Treat Claude Code as a development tool, not a demo: summarize source documents, generate structured reports, and turn a messy admin process into a repeatable workflow with review in the loop.
    \item Keep generated claims tied to source files so the agent cannot invent jobs, GPA, or Unreal experience. That is the same quality bar I would bring to internal tools that other developers have to trust.
\end{itemize}}}

\end{itemize}

% ============================================================
%     EDUCATION
% ============================================================

\section{Education}
\vspace{1pt}
\begin{itemize}

\needspace{5\baselineskip}
\item{\cventry{2022--2025}{Associate of Science in Computer Science}{Austin Community College}{Austin, TX}{}{\vspace{1pt}
Coursework foundation in programming and software development. LinkedIn-listed skills include C++, C\#, React, and game development. Seeking a paid internship that combines building developer tools with real production constraints.
}}

\end{itemize}

% ============================================================
%     TECHNICAL SKILLS
% ============================================================

\section{Technical Skills}
\vspace{1pt}
\begin{itemize}
\item \textbf{Languages:} TypeScript, JavaScript, SQL (SQLite), C++, C\#.
\item \textbf{Tools and platforms:} Node.js, discord.js, Git/GitHub, SQLite, REST APIs, Claude Code.
\item \textbf{Practices:} Cache-first API design, quota tracking, slash-command UX, embed reporting, README-level documentation, AI-assisted coding with human review.
\item \textbf{Not yet in public work:} Unreal Engine, Unreal Editor for Fortnite, C\# game features, desktop GUI frameworks, unit/functional test suites. Honest gaps. I learn unfamiliar tools on a live problem, which is how the Discord bot and Claude Code workflow got built.
\end{itemize}

% ============================================================
%     LANGUAGES
% ============================================================

\section{Languages}
\vspace{1pt}
\begin{itemize}
\item English.
\end{itemize}

% ============================================================
%     REFERENCES
% ============================================================

\section{References}
\vspace{1pt}
\begin{itemize}
\item{Available upon request.}
\end{itemize}

\end{document}
